% BÀI TOÁN HÌNH HỌC - BAI_HINH_37 (TRÍCH ĐỀ 37)
\begin{ex}
% Câu 17
	Cho đường tròn $(O)$, bán kính $R$, đường kính $AB$. Đường thẳng nằm ngoài đường tròn $d_1, d_2$ là các đường thẳng lần lượt qua $A, B$ và cùng vuông góc với đường kính $AB$. Qua điểm $H$ bất kì thuộc nửa đường tròn kẻ tiếp tuyến thứ ba cắt các tiếp tuyến qua $A, B$ tại $M, N$; $OM, ON$ cắt đường tròn lần lượt tại $E, F$. Gọi giao điểm của $ON$ và $HB$ là $G$; $OM$ và $HA$ là $K$, kẻ $HJ \perp AB$ ($J \in AB$).
	\begin{enumerate}[a)]
		\item Chứng minh tứ giác $AOHM$ nội tiếp.
		\item Chứng minh $JG^2 + JK^2 = R^2$. Xác định vị trí của $M, N$ để diện tích tam giác $MON$ đạt giá trị nhỏ nhất.
	\end{enumerate}
\loigiai{
{
}
\end{ex}
